\subsection{Large actual poles and fixed reciprocity do not force large indices}

\begin{proposition}[P-F33: actual period with zero strict indices]
\label{prop:largeperiodzeroindices}
For every integer $M\ge3$, give
$\mathbb C[x,y,z,w]$ the weights $(1,1,1,M)$ and set
\[
 R_M=\frac{\mathbb C[x,y,z,w]}
 {\,w(xy-z^2)+x^{M+2}+y^{M+2}\,}.
 \tag{P-F33a}
\]
This is a normal Cohen--Macaulay Gorenstein domain of
dimension three, with rational singularities and
$a$-invariant $-1$. Its Hilbert series is
\[
 H_M(t)=\frac{1-t^{M+2}}{(1-t)^3(1-t^M)},
 \qquad H_M(t^{-1})=-tH_M(t).
 \tag{P-F33b}
\]
It has a simple pole at every primitive $M$th root.
Nevertheless, if $f_M(s)=\dim(R_M)_s$, its global
quasipolynomial constituents of period $M$ are
\[
 P_r(X)=\left(\frac12+\frac1M\right)X(X+1)
                +\frac{(r+1)(M-r)}M,
 \qquad 0\le r<M.
 \tag{P-F33c}
\]
Their coefficients are all positive. In the strict
definitions of GCT VI,
\[
 p(f_M)=s(f_M)=0,
 \tag{P-F33d}
\]
while the minimum quasipolynomial period is exactly $M$.
Thus even actual poles at roots of arbitrarily large
multiplicative order, together with fixed
Krull dimension, normality, rational singularities,
Gorenstein symmetry, and a fixed negative $a$-invariant,
does not force an increasing positivity or saturation index.
No identification with an ordinary plethysm algebra is asserted.
\end{proposition}

\begin{proof}
Write $Q=xy-z^2$ and $F=x^{M+2}+y^{M+2}$.
The polynomial $Q$ is an irreducible quadratic in
$\mathbb C[x,y,z]$, and it does not divide $F$:
the point $(x,y,z)=(1,0,0)$ has $Q=0$ and $F=1$.
Hence $Q$ and $F$ are relatively prime. The polynomial
$wQ+F$ is primitive and linear in $w$ over this
unique factorization domain. It is irreducible over
its fraction field, so Gauss's lemma makes it
irreducible in the polynomial ring. It is therefore
prime, proving that $R_M$ is a domain of dimension three.
As a hypersurface it is Cohen--Macaulay and Gorenstein.

We check its singular locus rather than infer normality.
The partial derivatives of its equation are
\[
 wy+(M+2)x^{M+1},\quad
 wx+(M+2)y^{M+1},\quad -2wz,\quad xy-z^2.
 \tag{P-F33e}
\]
If $w\ne0$, their vanishing first gives $z=0$ and
$xy=0$. If $x=0$, the first derivative gives $y=0$;
if $y=0$, the second gives $x=0$.
If $w=0$, the first two derivatives give $x=y=0$,
and the fourth then gives $z=0$.
Conversely all four derivatives vanish on the entire
$w$-axis. The singular locus is exactly that axis.
It has codimension two, so the hypersurface is regular
in codimension one. Its Cohen--Macaulay property gives
$S_2$, and Serre's criterion proves normality.

The equation has weighted degree $M+2$ and the sum
of the four variable weights is $M+3$. The graded
hypersurface resolution, or equivalently its graded
$\operatorname{Ext}^1$ into the canonical module
of the polynomial ring, gives
\[
 \omega_{R_M}\simeq R_M((M+2)-(M+3))=R_M(-1).
 \tag{P-F33f}
\]
Thus $a(R_M)=-1$. The same one-step resolution
gives (P-F33b); direct substitution $t\mapsto t^{-1}$
verifies its displayed reciprocity identity.

Here are details verifying rational singularities.
Away from the origin, the only singular points
have $w\ne0$. On this open set, blow up the ideal
$(x,y,z)$. The $x$-chart of the strict transform,
with $y=xY$ and $z=xZ$, has equation
\[
 w(Y-Z^2)+x^M(1+Y^{M+2})=0.
 \tag{P-F33g}
\]
At its exceptional divisor $x=0$, the derivative
with respect to $Y$ is the unit $w$.
The $y$-chart is symmetric. In the $z$-chart,
$x=zX,\ y=zY$, the equation is
\[
 w(XY-1)+z^M(X^{M+2}+Y^{M+2})=0.
 \tag{P-F33h}
\]
At $z=0$, one has $XY=1$, so a derivative with
respect to $X$ or $Y$ is nonzero. Thus these charts
are smooth along the exceptional divisor, and
away from it the blowup is the identity on the
already smooth complement of the axis.
They give a resolution of $\operatorname{Spec}(R_M)_w$.

The exceptional divisor is the smooth conic
$xy-z^2=0$ in $\mathbb P^2\times\mathbb G_m$,
isomorphic to $\mathbb P^1\times\mathbb G_m$.
Its conormal line bundle is the restriction of
$\mathcal O_{\mathbb P^2}(1)$, hence is
the pullback of $\mathcal O_{\mathbb P^1}(2)$.
Writing $E$ for this divisor, the successive
quotients of its ideal powers are therefore
the pullbacks of $\mathcal O_{\mathbb P^1}(2j)$
for $j\ge0$.
Each has zero cohomology in positive degrees:
the base $\mathbb G_m$ is affine and
$H^i(\mathbb P^1,\mathcal O(2j))=0$ for $i>0$.
The exact sequences for successive thickenings
show the same vanishing for every infinitesimal
thickening of $E$.

By the theorem on formal functions, the completions
along the axis of all higher direct-image modules
of the structure sheaf under this resolution are
zero. These modules are coherent and supported on
the axis, since the resolution is an isomorphism
off it. A power of the axis ideal therefore
annihilates each module, so completion does not
change it. The modules themselves are zero.
This proves rationality on the punctured spectrum.

The graded rational-singularity criterion now applies:
Kurano--Sato--Singh--Watanabe,
\emph{Multigraded rings, diagonal subalgebras, and
rational singularities}, Theorem~5.6, printed p.~3262,
states that a normal $\mathbb N$-graded algebra of
finite type over a characteristic-zero field is
rational when it is Cohen--Macaulay, its punctured
spectrum is rational, and its $a$-invariant is negative.
All three hypotheses have just been checked.
Its formulation allows a local degree-zero ring;
here that ring is the field $\mathbb C$.
This proves the claimed rational singularities,
including at the graded origin.

We next compute every constituent and its strict index.
Write $s=Mk+r$ with $k\ge0$ and $0\le r<M$.
Using the power-series expansion of (P-F33b),
pair each summand after the first with the
corresponding shifted numerator summand. This gives
\begin{align*}
 f_M(s)
 &=\binom{s+2}{2}
      +\sum_{j=1}^{k}
           \left(\binom{s-jM+2}{2}
                        -\binom{s-jM}{2}\right)\\
 &=\binom{s+2}{2}
      +\sum_{j=1}^{k}\bigl(2(s-jM)+1\bigr)\\
 &=\binom{s+2}{2}+k(2s+1)-Mk(k+1).
 \tag{P-F33i}
\end{align*}
The pairing also holds for $r=0$ and $r=1$:
the formal expressions $\binom02$ and $\binom12$
are both zero, exactly matching the absent
negative-degree terms. Thus no initial exceptional
values have been omitted.
Substituting $k=(s-r)/M$ into (P-F33i) gives
exactly $P_r(s)$ in (P-F33c).
This proves global quasipolynomiality, including $s=0$.

The quadratic and linear coefficients of every
$P_r$ are $1/2+1/M>0$, and its constant coefficient
is $(r+1)(M-r)/M>0$. Hence all constituent coefficients
are nonnegative without any shift, which is the
definition of strict positivity in GCT VI,
Definition~1.2.2, printed p.~13. Moreover every
$P_r(X)$ is positive for every real $X\ge0$,
and in particular for every integer $X\ge1$.
This verifies strict saturation in Definition~1.2.4,
printed p.~14, on all constituent values, not only
on the residue class on which that constituent
represents $f_M$. Both indices are therefore zero.

Finally, if $\zeta$ is primitive of order $M$,
the numerator of (P-F33b) at $\zeta$ is
$1-\zeta^2\ne0$, since $M\ge3$.
The factor $1-t^M$ has a simple zero there
and $1-\zeta\ne0$. This is an actual simple pole.
Every quasipolynomial period must consequently
be divisible by $M$, by the elementary
generating-function argument used in P-F31.
Formula (P-F33c) already supplies period $M$,
so this is exactly the minimum period.
\end{proof}

\paragraph{Primary algebraic references and scope.}
Serre's criterion is Stacks Project
Lemma~10.157.4, Tag~031S; the complete-intersection
Gorenstein property is Lemma~47.21.8, Tag~0DWA.
The cohomology used for the exceptional conic is
Lemma~30.8.1, Tag~01XT, and the formal-functions
step is Theorem~30.20.5, Tag~02OC.
The rational-singularity criterion is available at
\url{https://www.math.utah.edu/~singh/publications/kurano_sato_singh_watanabe.pdf}.
The strict-index definitions are in
\url{https://arxiv.org/pdf/0704.0229}.

This proposition settles a side question about a
proposed inference from poles and reciprocity.
It neither bounds nor disproves PH2 or SH for the
ordinary cubic plethysm sequence in P-F31.
Any such conclusion needs additional
representation-specific information about its
constituent polynomials.

One further obstruction to an elementary shortcut is
visible in the same formula. The actual coefficients
satisfy $f_M(s)=\binom{s+2}{2}$ for every $0\le s<M$,
but the average of the constituent constant terms is
\[
 \frac1M\sum_{r=0}^{M-1}P_r(0)
       =\frac{(M+1)(M+2)}{6M}.
 \tag{P-F33j}
\]
To check the sum, expand
$(r+1)(M-r)=M(r+1)-r(r+1)$ and use
the finite sums of $r$ and $r^2$.
It grows linearly with $M$, despite the unchanged
initial coefficients. Thus those initial values alone
do not bound the off-residue constituent average
that would enter a Fourier comparison with a pole
coefficient.
